\section{Metrics and Discussion}

\begin{frame}{Why metrics deserve their own section}
  \SectionPill{Evaluation}{PresThreeMetric}
  \begin{columns}[T,onlytextwidth]
    \begin{column}{0.32\textwidth}
      \MiniCard{PresThreeMetric}{Effectiveness}{Does it work?\\[0.25em]{\scriptsize Acc \(\cdot\) F1 \(\cdot\) ASR \(\cdot\) CM \(\cdot\) ADR}}
    \end{column}
    \begin{column}{0.32\textwidth}
      \MiniCard{PresThreeMetric}{Efficiency}{What did it cost?\\[0.25em]{\scriptsize budget \(\Delta\) \(\cdot\) AML \(\cdot\) training cost}}
    \end{column}
    \begin{column}{0.32\textwidth}
      \MiniCard{PresThreeMetric}{Imperceptibility}{Is it noticeable?\\[0.25em]{\scriptsize similarity \(\cdot\) stats \(\cdot\) constraints}}
    \end{column}
  \end{columns}
  \vspace{0.45em}
  \begin{beamercolorbox}[wd=\linewidth,rounded=true,sep=0.55em]{takeaway box}
    {\color{PresThreeMetric}\bfseries Task \(\rightarrow\) metric map}\par
    \vspace{0.15em}
    \begin{tabular}{@{}p{0.34\linewidth}p{0.58\linewidth}@{}}
      classification & Acc \(\cdot\) F1 \(\cdot\) FNR \(\cdot\) FPR \\
      ranking / link prediction & AUC \(\cdot\) AP \(\cdot\) MRR \(\cdot\) Hits@K \\
      clustering / communities & NMI \(\cdot\) ARI \(\cdot\) modularity
    \end{tabular}
  \end{beamercolorbox}
  \SlideNote{2:00}{The paper treats metrics as more than bookkeeping. Metric choice shapes what robustness claims can mean, especially when effectiveness, efficiency, and imperceptibility point in different directions.}
\end{frame}

\begin{frame}{The few metrics worth remembering}
  \SectionPill{Evaluation formulas}{PresThreeMetric}
  \begin{columns}[T,onlytextwidth]
    \begin{column}{0.31\textwidth}
      \FormulaCard{PresThreeMetric}{Attack Success Rate}{
        \scriptsize
        \[
          ASR=\frac{\text{successful attacks}}{\text{attacks}}
        \]
      }
    \end{column}
    \begin{column}{0.31\textwidth}
      \FormulaCard{PresThreeMetric}{Average Modified Links}{
        \scriptsize
        \[
          AML=\frac{\text{modified links}}{\text{attacks}}
        \]
      }
    \end{column}
    \begin{column}{0.31\textwidth}
      \FormulaCard{PresThreeMetric}{Average Defense Rate}{
        \scriptsize
        \[
          ADR=\frac{ASR_{\text{with defense}}}{ASR_{\text{without defense}}}-1
        \]
      }
    \end{column}
  \end{columns}
  \vspace{0.45em}
  \begin{block}{Optional detail}
    \(CM\): margin to the correct class. \(\Lambda\): statistical similarity between clean and attacked graphs.
  \end{block}
  \SlideNote{2:30}{Keep only a few formulas. ASR is attack effectiveness, AML is perturbation cost, and ADR is the defense-side summary metric mentioned by the survey.}
\end{frame}

\begin{frame}{Open problems for attacks}
  \SectionPill{Discussion}{PresThreeAttack}
  \begin{columns}[T,onlytextwidth]
    \begin{column}{0.48\textwidth}
      \MiniCard{PresThreeAttack}{Unnoticeable attacks}{Effective perturbations should also preserve important graph properties.}
      \vspace{0.35em}
      \MiniCard{PresThreeAttack}{Scalable attacks}{Large graphs make full-gradient or full-candidate methods expensive.}
    \end{column}
    \begin{column}{0.48\textwidth}
      \MiniCard{PresThreeAttack}{Weaker knowledge}{Practical attacks should need less graph, label, and model access.}
      \vspace{0.35em}
      \MiniCard{PresThreeAttack}{Real-world graphs}{Move beyond simple, static, undistorted graph inputs.}
    \end{column}
  \end{columns}
  \SlideNote{2:00}{Phrase the frontier as the gap between benchmark attacks and realistic attacks. The difficult part is remaining effective while hidden, scalable, and practical.}
\end{frame}

\begin{frame}{Open problems for defenses and evaluation}
  \SectionPill{Discussion}{PresThreeMetric}
  \begin{columns}[T,onlytextwidth]
    \begin{column}{0.48\textwidth}
      \begin{exampleblock}{Defense}
        \begin{itemize}
          \item defense beyond node classification
          \item efficient defense at scale
          \item certification of robustness
        \end{itemize}
      \end{exampleblock}
    \end{column}
    \begin{column}{0.48\textwidth}
      \begin{block}{Evaluation}
        \begin{itemize}
          \item better cost metrics
          \item better imperceptibility metrics
          \item better metric selection by scenario
        \end{itemize}
      \end{block}
    \end{column}
  \end{columns}
  \vspace{0.35em}
  {\small\color{PresThreeMuted}Counting modified edges is often too crude: adding and removing an edge may have different real costs.}
  \SlideNote{2:00}{A useful message for a technical class is that evaluation itself remains underdeveloped. Graph imperceptibility is harder to define than an \(l_p\) constraint in images.}
\end{frame}

\begin{frame}{What the paper leaves us with}
  \SectionPill{Conclusion}{PresThreePrimary}
  \begin{block}{This survey's main contribution is structure}
    \begin{enumerate}
      \item Attacks are constrained perturbation problems on graphs.
      \item Defenses intervene at the data, model, objective, training, or monitoring level.
      \item Metrics and assumptions matter as much as raw performance.
    \end{enumerate}
  \end{block}
  \vspace{0.45em}
  \centering
  \begin{tikzpicture}
    \node[modelbox,draw=PresThreeAttack,fill=PresThreeAttackSoft,text width=3.2cm] (a) {unified attack};
    \node[font=\large\bfseries,right=0.8cm of a] (arrow) {\(\leftrightarrow\)};
    \node[modelbox,draw=PresThreeDefense,fill=PresThreeDefenseSoft,text width=3.2cm,right=0.8cm of arrow] (d) {unified defense};
  \end{tikzpicture}
  \SlideNote{1:45}{The main contribution is not one best attack or one best defense. It is a structured language for thinking about graph adversarial learning.}
\end{frame}

\begin{frame}[plain]
  \vspace{0.75cm}
  \begin{columns}[c,onlytextwidth]
    \begin{column}{0.62\textwidth}
      {\fontsize{24}{28}\selectfont\bfseries A Survey of Adversarial Learning on Graphs\par}
      \vspace{0.35cm}
      {\fontsize{13}{16}\selectfont\color{PresThreeMuted}Reliable graph learning needs better models and better threat models.\par}
      \vspace{0.45cm}
      {\small Krzysztof Czerwionka and Przemysław Fuchs}
    \end{column}
    \begin{column}{0.34\textwidth}
      \centering
      \ClosingGraph[1.05]
    \end{column}
  \end{columns}
  \SlideNote{0:45}{Clean landing slide. Leave it up for questions or for the natural end of the talk.}
\end{frame}
